\documentclass{article}
\usepackage{amsmath, amssymb, amsthm}
\usepackage{geometry}
\usepackage[UTF8]{ctex}

\newtheorem{problem}{问题}
\newtheorem{solution}{解答}
\title{数值分析第八周作业}
\date{\today}
\author{李佩优PB23000270}
\geometry{a4paper,scale=0.8}

\begin{document}
\maketitle

\section*{7.4}
\subsection*{8}

% 解： 应用龙贝格算法，先计算梯形法则的序列 \(T_1, T_2, T_4\)，再通过理查德森外推得到 \(R(2,2)\)。

\subsubsection*{a}  
\(\displaystyle \int_{1}^{3} \frac{dx}{x}\)

区间 \([1,3]\)，长度 \(b-a=2\)。

\(T_1\)（1个区间）：
  \[
  T_1 = \frac{2}{2}\left(f(1)+f(3)\right) = 1\times\left(1+\frac{1}{3}\right) = \frac{4}{3}.
  \]

\(T_2\)（2个区间，步长 \(h=1\)）：
  \[
  T_2 = \frac{1}{2}\left(f(1)+2f(2)+f(3)\right) = \frac{1}{2}\left(1+2\times\frac{1}{2}+\frac{1}{3}\right) = \frac{1}{2}\times\frac{7}{3} = \frac{7}{6}.
  \]

\(T_4\)（4个区间，步长 \(h=0.5\)）：
  \[
  T_4 = \frac{0.5}{2}\left(f(1)+2f(1.5)+2f(2)+2f(2.5)+f(3)\right) = 0.25\left(1+2\times\frac{2}{3}+2\times\frac{1}{2}+2\times\frac{2}{5}+\frac{1}{3}\right) = \frac{67}{60}.
  \]

构造龙贝格表：
\[
R(0,0)=T_1=\frac{4}{3},\quad R(1,0)=T_2=\frac{7}{6},\quad R(2,0)=T_4=\frac{67}{60}.
\]
外推：
\[
R(1,1)=\frac{4R(1,0)-R(0,0)}{3}=\frac{4\cdot\frac{7}{6}-\frac{4}{3}}{3}=\frac{10}{9},
\]
\[
R(2,1)=\frac{4R(2,0)-R(1,0)}{3}=\frac{4\cdot\frac{67}{60}-\frac{7}{6}}{3}=\frac{11}{10},
\]
\[
R(2,2)=\frac{16R(2,1)-R(1,1)}{15}=\frac{16\cdot\frac{11}{10}-\frac{10}{9}}{15}=\frac{742}{675}.
\]

故 \(R(2,2)=\dfrac{742}{675}\)。

\subsubsection*{b}  
\(\displaystyle \int_{0}^{\pi/2} \left(\frac{x}{\pi}\right)^2 dx\)

区间 \([0,\pi/2]\)，长度 \(b-a=\pi/2\)，被积函数 \(f(x)=x^2/\pi^2\)。

\(T_1\)：
  \[
  T_1 = \frac{\pi/2}{2}\left(f(0)+f(\pi/2)\right) = \frac{\pi}{4}\left(0+\frac{1}{4}\right) = \frac{\pi}{16}.
  \]

\(T_2\)（步长 \(h=\pi/4\)）：
  \[
  T_2 = \frac{\pi/4}{2}\left(f(0)+2f(\pi/4)+f(\pi/2)\right) = \frac{\pi}{8}\left(0+2\cdot\frac{1}{16}+\frac{1}{4}\right) = \frac{\pi}{8}\cdot\frac{3}{8} = \frac{3\pi}{64}.
  \]

\(T_4\)（步长 \(h=\pi/8\)）：
  \[
  T_4 = \frac{\pi/8}{2}\left(f(0)+2f(\pi/8)+2f(\pi/4)+2f(3\pi/8)+f(\pi/2)\right) = \frac{\pi}{16}\left(0+2\cdot\frac{1}{64}+2\cdot\frac{1}{16}+2\cdot\frac{9}{64}+\frac{1}{4}\right) = \frac{11\pi}{256}.
  \]

龙贝格表：
\[
R(0,0)=\frac{\pi}{16},\quad R(1,0)=\frac{3\pi}{64},\quad R(2,0)=\frac{11\pi}{256}.
\]
外推：
\[
R(1,1)=\frac{4\cdot\frac{3\pi}{64}-\frac{\pi}{16}}{3}=\frac{\frac{12\pi}{64}-\frac{4\pi}{64}}{3}=\frac{\frac{8\pi}{64}}{3}=\frac{\pi}{24},
\]
\[
R(2,1)=\frac{4\cdot\frac{11\pi}{256}-\frac{3\pi}{64}}{3}=\frac{\frac{44\pi}{256}-\frac{12\pi}{256}}{3}=\frac{\frac{32\pi}{256}}{3}=\frac{\pi}{24},
\]
\[
R(2,2)=\frac{16\cdot\frac{\pi}{24}-\frac{\pi}{24}}{15}=\frac{\frac{15\pi}{24}}{15}=\frac{\pi}{24}.
\]

故 \(R(2,2)=\dfrac{\pi}{24}\)。

\end{document}